\scanpage{56}
% manuscript line(s) 1–4
des couples $(n,\unsure{\mathscr X|\mathcal X|\mathfrak X})$ où $n\in\mathbb Z/2\mathbb Z$ est une ``parité'', $\mathscr X/X$ rev. quadratique \uline{avec} \quad $X'\in\underline{\mathrm{Rev.quad}}^{!}$ \uline{si $n=1$}

% manuscript line(s) 5
On pose
% manuscript line(s) 6–7
\[
(n,\deleted{\unsure{X'}}\mathscr X)\wedge(n',\mathscr X')=\begin{cases}(n+n',\mathscr X\wedge\mathscr X')&\text{si $n$ ou $n'$ pair}\\ (n+n',\mathscr X\overset{!}{\wedge}\mathscr X')&\text{si $n$ \uline{et} $n'$ impairs}\end{cases}
\]
% manuscript line(s) 8–10
où $\mathscr X\overset{!}{\wedge}\mathscr X'$ est un produit contracté modifié, défini si $\mathscr X,\mathscr X'\in\underline{\mathrm{Rev.quad}}^!(X)$, et défini (si $2$ rég. sur $X$) par la condition
% manuscript line(s) 11
\[
(L,\delta,0)\overset{!}{\wedge}(L',\delta',0)=(L\otimes L',\tfrac14\delta\delta',0)
\]
% manuscript line(s) 12–13
où $\tfrac14\delta\delta'$ a un sens intrinsèque, car loc.\textsuperscript{t} on a
% manuscript line(s) 13–14
\[
\begin{aligned}\delta&=4\Delta &&\Delta\ \text{définie mod }h\ \text{avec }4h=0\\ \delta'&=4\Delta' &&\Delta'\ \text{\rule{1.2cm}{.3pt} }h'\ \text{\rule{1.2cm}{.3pt} }4h'=0\end{aligned}
\]
% manuscript line(s) 15
\[
\delta\delta'=16\Delta\Delta'\qquad\tfrac14\delta\delta'\overset{\mathrm{def}}=4\Delta\Delta'
\]
% manuscript line(s) 15–16
qui ne change pas \unsure{--} si $\Delta\longmapsto\Delta+h\quad(4h=0$ \quad $\Delta'\longmapsto\Delta'+h'\quad(4h'=0)$

% manuscript line(s) 17–19
\leavevmode\llap{\rule[-2\baselineskip]{.3pt}{3\baselineskip}\hspace{2pt}\rule[-2\baselineskip]{.3pt}{3\baselineskip}\quad}Il faut lui donner un sens indépendamment de la condition $2$ régulier sur $X$, exercice catégorique universel $\ldots$

% manuscript line(s) 20–24
Cette opération est \uline{commutative} \inserted{(à isom. can. près)} et de plus associative, les isom de comm. et d'assoc. étant les habituels sur $L\otimes L',\deleted{\unsure{\simeq}}L'\otimes L$, et $(L\otimes L')\otimes L''\simeq L\otimes(L'\otimes{\unsure{\underline L|L}}'')$, et respectant automatiquement les pondérations $\delta,T_0$ etc.
% archive sheet number: 55
